% ============================================================================
% AD7031 Homework 3 -- covers Weeks 9-13 (chance constraints, classification
% as a MILP with gurobipy, risk measures VaR and CVaR in the objective
% and in a constraint).
% 80 points, 4 problems of 20.
% Problem 1 starts from Week 10's chance constraint program: its scenario MIP
% and the OR-vs-ML notation table. Problem 2 is the coding problem:
% classification in the slide deck's theta/(x,y) notation as the training
% MILP (OL7014 midterm Problem 2), with a gurobipy skeleton whose blanks are
% marked in red. Problems 3-4 grow out of the risk-measures problem dropped
% from HW2 (not yet covered then): Problem 3 is VaR, Problem 4 is CVaR, and
% each one asks the same five things -- meaning, by-hand values at three
% epsilons, reformulation in the objective, reformulation in a constraint,
% the full (SP). The two-stage LAP problem was dropped (Nov 14).
% Self-contained on purpose -- no \input{preamble} -- compiles with pdflatex.
% ============================================================================
\documentclass[10pt]{article}

\usepackage[a4paper, margin=2.2cm]{geometry}
\usepackage{amsmath, amssymb}
\usepackage[dvipsnames]{xcolor}
\usepackage{array}
\usepackage{enumitem}
\usepackage{fancyvrb}
\usepackage[hidelinks]{hyperref}

\definecolor{ruleink}{RGB}{18, 68, 140}
\definecolor{keyink}{RGB}{20, 90, 20}
\definecolor{codeink}{RGB}{190, 30, 30}

% ---- "write your code here" marker inside a Verbatim block ------------------
\newcommand{\yourcode}[1]{\textcolor{codeink}{\textbf{#1}}}

\setlength{\parindent}{0pt}
\pagestyle{plain}

\AtBeginDocument{%
  \abovedisplayskip=5pt plus 1pt
  \belowdisplayskip=5pt plus 1pt
  \abovedisplayshortskip=3pt
  \belowdisplayshortskip=3pt}

% ---- one problem -----------------------------------------------------------
\newcommand{\problem}[2]{%   points, title
  \par\bigskip
  {\color{ruleink}\rule{\textwidth}{1pt}}\par\nobreak\vspace{3pt}
  {\large\bfseries #2}\hfill{\normalsize\bfseries(#1 points)}\par\nobreak\vspace{4pt}}

% ---- one-line hint, blue italic --------------------------------------------
\newcommand{\hint}[2][]{%   optional number: \hint[1]{...} -> "hint 1:"
  \par\vspace{3pt}{\color{ruleink}\emph{hint\if\relax\detokenize{#1}\relax\else\ #1\fi:\ #2}}\par}

\newlist{parts}{enumerate}{1}
\setlist[parts]{label=(\alph*), leftmargin=2.2em, itemsep=5pt}

\newcommand{\VaR}{\mathrm{VaR}}
\newcommand{\CVaR}{\mathrm{CVaR}}
\newcommand{\E}{\mathbb{E}}
\newcommand{\Prob}{\mathbb{P}}

\begin{document}

\begin{center}
  {\large AD7031 Optimization Under Uncertainty}\\[3pt]
  {\Large\bfseries Homework 3}\\[4pt]
  {\small \textbf{Due Nov 28, start of class}}
\end{center}
\medskip\hrule\medskip

Write down your solutions and submit them as a single PDF to
\texttt{lossof0706@gmail.com}.\\
\textbf{\underline{Your solutions must be handwritten.}}

% ================================================================ Problem 1
\problem{20}{Problem 1. Chance constraints}

Recall the chance-constrained program of Weeks 10--11:
\begin{equation}
\label{eq:ccp}
\text{(CCP)} \quad
\begin{aligned}
  \min_{x} \quad & f(x) \\
  \text{s.t.} \quad & x \in \mathcal{X} \\
                    & \Prob\big[\, g(x, \tilde\xi) \le 0 \,\big] \;\ge\; 1 - \varepsilon ,
\end{aligned}
\end{equation}
where $x$ is the decision, $\mathcal{X}$ collects the constraints with no uncertainty, $\tilde\xi \sim \mathbb{P}$ is the
uncertain data, $g$ is the one constraint that involves it, and $\varepsilon \in (0,1)$ is a tolerance.
Throughout this homework, ``$S$ scenarios'' means $\tilde\xi$ takes the values $\xi^1, \dots, \xi^S$ with
probabilities $p_1, \dots, p_S$; with data, $p_s = \tfrac1S$ (the empirical distribution of $S$ data points, as in Week 10).

\begin{parts}
\item \textbf{(6 pts)} Explain in plain Korean what the constraint $\Prob[\, g(x, \tilde\xi) \le 0 \,] \ge 1 - \varepsilon$
  means. Give one potential example of a chance constraint in military applications.

\item \textbf{(8 pts)} With $S$ scenarios, reformulate \eqref{eq:ccp} as a mixed-integer program,
  with one binary $z_s$ per scenario. \textbf{Show the derivation clearly.}
  \hint{$z_s = 1$ must force $g(x, \xi^s) \le 0$; $z_s = 0$ must leave the row saying nothing.}

\item \textbf{(6 pts)} \textbf{One problem, two languages.} The slide deck of Week 10 writes the same program in the
  language of machine learning (ML). Fill in the ML column:
  \begin{center}
  \begin{tabular}{l|c|c}
     & OR & ML \\ \hline
    decision variable & $x$ & \rule{2.5cm}{0.4pt} \\
    uncertain data & $\tilde\xi$ & $\tilde\xi = ($\rule{1.2cm}{0.4pt}$,\ $\rule{1.2cm}{0.4pt}$)$ \\
    model & --- & $f_{\theta}(x)$ \\
    objective function & $f(x, \tilde\xi)$ & \rule{2.5cm}{0.4pt} \\
  \end{tabular}
  \end{center}
\end{parts}

% ================================================================ Problem 2
\problem{20}{Problem 2. Classification as a MILP}

\textbf{Classification, in the ML language.} $N = 5$ samples $\xi_i = (x_i, y_i)$: an input
$x_i \in \mathbb{R}^3$ whose last entry is the constant $1$ (the bias term), and a label $y_i \in \{-1, +1\}$:
\[
\begin{array}{c|cc}
  i & x_i^{\top} & y_i \\ \hline
  1 & (1,\,2,\,1) & +1 \\
  2 & (2,\,3,\,1) & +1 \\
  3 & (3,\,1,\,1) & +1 \\
  4 & (0,\,0,\,1) & -1 \\
  5 & (2,\,2,\,1) & -1
\end{array}
\]
The linear model $f_{\theta}(x) = \theta^{\top}x$, $\theta \in \mathbb{R}^3$, predicts $+1$ when $\theta^{\top}x > 0$ and
$-1$ otherwise. As in Week 10, sample $i$ is classified correctly \emph{with margin} when
$y_i\,\theta^{\top}x_i \ge 1$, i.e., when $g(\theta, \xi_i) := 1 - y_i\,\theta^{\top}x_i \le 0$. Use $M = 10$.

\begin{parts}
\item \textbf{(10 pts)} Write the binary classification problem as a mixed-integer program. For $x_i$ and $y_i$,
  use the numbers given in the table.
  \hint{OL7014's midterm Problem 2, with the margin: $z_i = 1$ forces $1 - y_i\,\theta^{\top}x_i \le 0$.}

\item \textbf{(10 pts)} Implement (a) with \texttt{gurobipy}
  and \textbf{paste a screenshot of your code and of its printed output}.
  Following is skeleton code with the data typed in; write your code where it says so, in
  \textcolor{codeink}{\textbf{red}}:
\begin{Verbatim}[commandchars=\\\{\}, fontsize=\small]
import gurobipy as gp
from gurobipy import GRB

X = [[1, 2, 1], [2, 3, 1], [3, 1, 1], [0, 0, 1], [2, 2, 1]]   # inputs x_i (bias 1 last)
y = [+1, +1, +1, -1, -1]                                       # labels y_i
N, M = 5, 10

model = gp.Model("classifier")           # on Colab: gp.Model("classifier", env=env)

# variables
theta = \yourcode{...   # three free continuous variables}
z     = \yourcode{...   # one binary per sample}

# constraints
\yourcode{...   # one big-M row per sample}

# objective
\yourcode{...}

model.optimize()
print(f"optimal theta is \{[theta[j].X for j in range(3)]\}")
print(f"z is \{[int(round(z[i].X)) for i in range(N)]\}")
print(f"optimal objective value is \{model.ObjVal\}")
\end{Verbatim}
\end{parts}

% ================================================================ Problem 3
\problem{20}{Problem 3. Value at Risk}

Recall the general stochastic program of Week 11,
\[
\text{(SP)} \quad
\begin{aligned}
  \min_{x} \quad & \mathcal{R}_{\mathbb{P}}\big[f(x, \tilde\xi)\big] \\
  \text{s.t.} \quad & x \in \mathcal{X} \\
                    & \mathcal{R}_{\mathbb{P}}\big[g(x, \tilde\xi)\big] \le 0 ,
\end{aligned}
\]
where $\tilde\xi \sim \mathbb{P}$ and $\mathcal{R}_{\mathbb{P}}$ is a risk measure of our choice. Problems 3 and 4 take
the two risk measures of the course in turn, and ask the same questions of each: what it means, how to compute it,
how to reformulate it \emph{in the objective}, and how to reformulate it \emph{in a constraint}.

\begin{parts}
\item \textbf{(3 pts)} Suppose we choose $\mathcal{R}_{\mathbb{P}} = \VaR_{\varepsilon}$, i.e., the objective function in (SP) is
  \[
  \begin{aligned}
    \VaR_{\varepsilon}\big[f(x, \tilde\xi)\big] \;=\; \min_{t \in \mathbb{R}} \quad & t \\
    \text{s.t.} \quad & \Prob\big(f(x, \tilde\xi) \le t\big) \ge 1 - \varepsilon .
  \end{aligned}
  \]
  Using this minimization problem, explain in plain Korean what $\VaR_{\varepsilon}\big[f(x, \tilde\xi)\big]$ means.

\item \textbf{(3 pts)} Fix your decision at $x = \bar x$ and write $\tilde Z = f(\bar x, \tilde\xi)$ for the
  random objective value --- a random \emph{loss} (negative means a gain). Suppose $\tilde\xi$ takes only ten
  scenarios, so $\tilde Z$ takes the ten values below with the given probabilities:
  \begin{center}
  \begin{tabular}{c|cccccccccc}
    scenario $s$ & 1 & 2 & 3 & 4 & 5 & 6 & 7 & 8 & 9 & 10 \\ \hline
    $Z^s = f(\bar x, \xi^s)$ & $-12$ & $-8$ & $-5$ & $-3$ & $-1$ & $0$ & $2$ & $5$ & $9$ & $18$ \\
    $p_s = \Prob(\tilde Z = Z^s)$ & $0.1$ & $0.1$ & $0.1$ & $0.1$ & $0.1$ & $0.1$ & $0.1$ & $0.1$ & $0.1$ & $0.1$ \\
  \end{tabular}
  \end{center}
  \textbf{By hand}, what is $\VaR_{\varepsilon}[\tilde Z]$ when
  \begin{itemize}[leftmargin=2.2em, itemsep=2pt, topsep=3pt]
    \item $\varepsilon = 0.2$;
    \item $\varepsilon = 0.5$;
    \item $\varepsilon = 1.0$?
  \end{itemize}
  \hint{the $\VaR$ is the smallest value $t$ with $\Prob(\tilde Z \le t) \ge 1 - \varepsilon$; for $\varepsilon = 1.0$,
    ask what the constraint still requires.}

\item \textbf{(5 pts)} \textbf{VaR in the objective.} Recall
  \begin{equation}
  \label{eq:var-obj}
    \min_{x \in \mathcal{X}}\ \VaR_{\varepsilon}\big[f(x, \tilde\xi)\big]
    \qquad\Longleftrightarrow\qquad
    \begin{aligned}
      \min_{x \in \mathcal{X},\ t \in \mathbb{R}} \quad & t \\
      \text{s.t.} \quad & \Prob\big(f(x, \tilde\xi) \le t\big) \ge 1 - \varepsilon .
    \end{aligned}
  \end{equation}
  With $S$ scenarios, reformulate \eqref{eq:var-obj} as a mixed-integer program.
  \textbf{Show the derivation clearly.}
  \hint{the right-hand side of \eqref{eq:var-obj} is a CCP: Problem 1(b) with $f(x, \xi) - t$ in the role of $g$ ---
    only now $t$ is a decision variable too.}

\item \textbf{(5 pts)} \textbf{VaR in a constraint.} (For (d) and (e), you do not need to use the numbers in the
  table.) Show the reformulation of the constraint
  $\VaR_{\varepsilon}\big[g(x, \tilde\xi)\big] \le 0$, and conclude that it is \emph{exactly} the chance constraint of
  Problem 1, $\Prob\big(g(x, \tilde\xi) \le 0\big) \ge 1 - \varepsilon$. Then, with $S$ scenarios, write it as big-$M$
  rows with binaries $z_s$.
  \hint{use
    \[
      \VaR_{\varepsilon}\big[g(x, \tilde\xi)\big] \le 0
      \quad \Longleftrightarrow \quad
      \left[\begin{aligned}
        \min_{t \in \mathbb{R}} \quad & t \\
        \textnormal{s.t.} \quad & \Prob\big(g(x, \tilde\xi) \le t\big) \ge 1 - \varepsilon
      \end{aligned}\right] \le 0
      \quad \Longleftrightarrow \quad
      \begin{aligned}
        & t \le 0 \\
        & \Prob\big(g(x, \tilde\xi) \le t\big) \ge 1 - \varepsilon .
      \end{aligned}
    \]}

\item \textbf{(4 pts)} \textbf{Putting it together.} Write (SP) with $\mathcal{R}_{\mathbb{P}} = \VaR_{\varepsilon}$ and $S$
  scenarios as one mixed-integer program, using (c) for the objective and (d) for the constraint.
  How many binaries does it have? Explain in plain Korean why this problem is hard to solve even when $f$ and $g$
  are convex in $x$.
  \hint{the objective and the constraint each need their own binaries --- call them $z^f_s$ and $z^g_s$.}
\end{parts}

% ================================================================ Problem 4
\problem{20}{Problem 4. Conditional Value at Risk}

Same (SP) and the same ten-scenario table as Problem 3, now with $\mathcal{R}_{\mathbb{P}} = \CVaR_{\varepsilon}$.

\begin{parts}
\item \textbf{(3 pts)} The objective function in (SP) is now
  \[
    \CVaR_{\varepsilon}\big[f(x, \tilde\xi)\big] \;=\; \min_{\beta \in \mathbb{R}}\
      \beta + \frac{1}{\varepsilon}\, \E_{\mathbb{P}}\big[\max\{f(x, \tilde\xi) - \beta,\ 0\}\big] .
  \]
  Using this minimization problem, explain in plain Korean what $\CVaR_{\varepsilon}\big[f(x, \tilde\xi)\big]$ means.

\item \textbf{(3 pts)} The same ten scenarios as Problem 3(b), $\tilde Z = f(\bar x, \tilde\xi)$:
  \begin{center}
  \begin{tabular}{c|cccccccccc}
    scenario $s$ & 1 & 2 & 3 & 4 & 5 & 6 & 7 & 8 & 9 & 10 \\ \hline
    $Z^s = f(\bar x, \xi^s)$ & $-12$ & $-8$ & $-5$ & $-3$ & $-1$ & $0$ & $2$ & $5$ & $9$ & $18$ \\
    $p_s = \Prob(\tilde Z = Z^s)$ & $0.1$ & $0.1$ & $0.1$ & $0.1$ & $0.1$ & $0.1$ & $0.1$ & $0.1$ & $0.1$ & $0.1$ \\
  \end{tabular}
  \end{center}
  \textbf{By hand}, what is $\CVaR_{\varepsilon}[\tilde Z]$ when
  \begin{itemize}[leftmargin=2.2em, itemsep=2pt, topsep=3pt]
    \item $\varepsilon = 0.2$;
    \item $\varepsilon = 0.5$;
    \item $\varepsilon = 1.0$?
  \end{itemize}
  For each $\varepsilon$, compare $\CVaR_{\varepsilon}[\tilde Z]$ with $\VaR_{\varepsilon}[\tilde Z]$ from Problem 3(b), and
  explain in plain Korean the relation between the two.
  \hint{evaluate $\beta + \frac{1}{\varepsilon}\E[\max\{\tilde Z - \beta, 0\}]$ at $\beta = \VaR_{\varepsilon}[\tilde Z]$
    from Problem 3(b) --- that is where the minimum is attained; for $\varepsilon = 1.0$, compare with $\E[\tilde Z]$.}

\item \textbf{(5 pts)} \textbf{CVaR in the objective.} Recall
  \begin{equation}
  \label{eq:cvar-obj}
    \min_{x \in \mathcal{X}}\ \CVaR_{\varepsilon}\big[f(x, \tilde\xi)\big]
    \qquad\Longleftrightarrow\qquad
    \min_{x \in \mathcal{X},\ \beta \in \mathbb{R}}\ \
      \beta + \frac{1}{\varepsilon}\, \E_{\mathbb{P}}\big[\max\{f(x, \tilde\xi) - \beta,\ 0\}\big] .
  \end{equation}
  With $S$ scenarios, reformulate \eqref{eq:cvar-obj} with one variable $u_s \ge 0$ per scenario in place of each
  $\max\{\cdot, 0\}$ --- no binaries. \textbf{Show the derivation clearly.}
  \hint{$u_s \ge f(x, \xi^s) - \beta$ and $u_s \ge 0$ replace $\max\{f(x, \xi^s) - \beta,\ 0\}$ --- why is that
    allowed in a minimization?}

\item \textbf{(5 pts)} \textbf{CVaR in a constraint.} Show the reformulation of the constraint
  $\CVaR_{\varepsilon}\big[g(x, \tilde\xi)\big] \le 0$ with its own variable $\beta_g$. Then, with $S$ scenarios,
  write it as rows with variables $u^g_s \ge 0$.
  \hint{do the same as in Problem 3(d): $[\min_{\beta_g} \cdots] \le 0 \iff \exists\, \beta_g$ with $\cdots \le 0$.}

\item \textbf{(4 pts)} \textbf{Putting it together.} Write (SP) with $\mathcal{R}_{\mathbb{P}} = \CVaR_{\varepsilon}$ and $S$
  scenarios as one optimization problem, using (c) for the objective and (d) for the constraint.
  Compare with Problem 3(e): is this problem convex when $f$ and $g$ are convex in $x$? Explain in plain Korean
  why, for computational tractability, we may prefer $\CVaR$ over $\VaR$.
  \hint{one $\beta$ per risk term: $\beta_f$ for the objective and $\beta_g$ for the constraint; both join the outer
    $\min$, with $u^f_s$ and $u^g_s$.}
\end{parts}

\end{document}
